% RuSMT — SPECops 2026 submission.
% ACM SIGPLAN acmart, two-column, DOUBLE-BLIND (sigplan,review,anonymous).
% De-anonymize for camera-ready by removing the `anonymous' class option and
% filling in the author block flagged with %DEANON below.
\documentclass[sigplan,review,anonymous]{acmart}

% --- workshop / front matter (no DOI for a workshop submission) ---
\settopmatter{printacmref=false}
\setcopyright{none}
\renewcommand\footnotetextcopyrightpermission[1]{}
\acmConference[SPECops 2026]{1st International Workshop on Specification-Driven
  Development Life Cycle}{October 2026}{co-located with ISSTA, Oakland, CA, USA}
\acmYear{2026}
\settopmatter{printfolios=true}

% --- packages on top of acmart ---
\usepackage{booktabs}
\usepackage{listings}
\usepackage{tikz}
\usetikzlibrary{arrows.meta,positioning,calc,shapes.geometric,fit}
\usepackage{enumitem}

% acmart loads amsmath/amsthm and defines theorem environments; add our own.
\theoremstyle{definition}
\newtheorem{principle}{Principle}
\theoremstyle{remark}
\newtheorem{observation}{Observation}
% acmart already provides `definition'.

\newcommand{\rusmt}{RuSMT}
\newcommand{\imp}{IMP}
\newcommand{\toml}{TOML}
\newcommand{\code}[1]{\texttt{#1}}
\setlength{\emergencystretch}{3em}

\definecolor{cmtcol}{gray}{0.4}
\lstdefinestyle{rust}{
  basicstyle=\ttfamily\scriptsize,
  keywordstyle=\bfseries,
  commentstyle=\itshape\color{cmtcol},
  stringstyle=\color{teal},
  morekeywords={fn,let,match,if,else,enum,struct,pub,return,use,mod,as,for,in},
  morecomment=[l]{//},
  morestring=[b]",
  columns=fullflexible,
  keepspaces=true,
  showstringspaces=false,
  breaklines=true,
  captionpos=b,
  frame=tb,
  framerule=0.3pt,
  xleftmargin=0.4em,
  aboveskip=0.5em,belowskip=0.3em,
}
\lstdefinestyle{smt}{
  basicstyle=\ttfamily\scriptsize,
  columns=fullflexible,
  keepspaces=true,
  showstringspaces=false,
  breaklines=true,
  captionpos=b,
  frame=tb,
  framerule=0.3pt,
  xleftmargin=0.4em,
  aboveskip=0.5em,belowskip=0.3em,
}
\lstset{style=rust}

\newcommand{\certcount}{131}
\newcommand{\misscount}{51}
\newcommand{\proposalcount}{689}
\newcommand{\rejectcount}{551}
\newcommand{\undecidedcount}{7}
\newcommand{\diffcount}{15}
\newcommand{\markers}{182}
\newcommand{\impmarkers}{2}

\begin{document}

\title{\rusmt{}: An Executable Semantics as Conformance Oracle and Test Suite Synthesizer}

%DEANON: replace the anonymous block below at camera-ready with real authors.
\author{Anonymous Author(s)}
\affiliation{%
  \institution{Anonymous Institution}
  \country{}
}


\begin{abstract}
Conformance suites are most valuable where specifications are narrow and
negative: malformed syntax, duplicate declarations, out-of-range literals, and
dynamic errors. These cases are easy to name in a reference semantics and tedious
to construct by hand. We present \rusmt{}, a framework that turns one executable
semantics into both a conformance oracle and a marker-directed test-suite
synthesizer. The author writes the semantics as a restricted Rust program over
SMT-denoting types and marks obligations with stable names such as
\code{Path::named("division\_by\_zero")}. The same source runs concretely as the
oracle and lowers to one SMT-LIB reachability query per marker; each accepted
model becomes a test input.

The main engineering problem is not expressing the semantics but finding inputs.
On a \toml{}~1.1.0 parser with \markers{} markers, Z3 alone produced no witness
within affordable budgets: symbolic parser inputs keep both sides of deeply
nested branches alive during macro expansion. \rusmt{} therefore uses a two-stage
pipeline. Stage~1 asks Z3 to synthesize the input directly. Stage~2, used only
after timeout or \code{unknown}, asks a language model to propose a candidate
from the emitted SMT-LIB query and Z3's previous responses. That candidate is
not trusted; it is added as an equality to the original query, and only a
\code{sat} result from Z3 becomes a witness. This converts the hard part of the
\toml{} run from solver search to solver acceptance while preserving the original
marker obligation.

In our run, Stage~1 found both \imp{} witnesses (\impmarkers{}/\impmarkers{}) and
no \toml{} witnesses (0/\markers{}). Stage~2, using live \code{gpt-5.4-mini}
calls and cached transcripts for replay, produced \certcount{}/\markers{}
\toml{} witnesses; \misscount{} markers remain uncovered and are reported by
name. The same run recorded \proposalcount{} candidate proposals: \certcount{}
were accepted, \rejectcount{} were rejected by Z3 as not reaching their claimed
marker, and \undecidedcount{} were undecided. Running the generated \toml{} suite
against four independent parsers produced \diffcount{} accept/reject
divergences, including likely parser bugs and several float-boundary cases that
require specification triage. The contribution is the pipeline and artifact: an
executable semantics whose marked branches generate a reusable conformance suite,
with Z3 providing logical acceptance and replay checking the SMT lift against the
concrete oracle.
\end{abstract}

\keywords{executable semantics, conformance testing, bounded program synthesis,
SMT, Z3, solver-aided programming, Rust, large language models}

\begin{CCSXML}
<ccs2012>
<concept>
<concept_id>10011007.10011074.10011099.10011102.10011103</concept_id>
<concept_desc>Software and its engineering~Software testing and debugging</concept_desc>
<concept_significance>500</concept_significance>
</concept>
<concept>
<concept_id>10011007.10011074.10011099.10011692</concept_id>
<concept_desc>Software and its engineering~Formal software verification</concept_desc>
<concept_significance>500</concept_significance>
</concept>
<concept>
<concept_id>10003752.10003790.10011740</concept_id>
<concept_desc>Theory of computation~Automated reasoning</concept_desc>
<concept_significance>300</concept_significance>
</concept>
</ccs2012>
\end{CCSXML}

\ccsdesc[500]{Software and its engineering~Software testing and debugging}
\ccsdesc[500]{Software and its engineering~Formal software verification}
\ccsdesc[300]{Theory of computation~Automated reasoning}

\maketitle

% =====================================================================
\section{Introduction}
\label{sec:intro}

A language ecosystem typically carries three artifacts meant to say the same
thing: a prose \emph{specification}, an \emph{implementation} adhering to the specification, and a
\emph{test generator}. Written separately, they drift. Fuzzing covers the bulk
of the input space but misses the narrow error and edge conditions a
specification cares about most; a hand-written suite forces its author both to
decide which conditions are interesting and to craft inputs exercising them.
\rusmt{} keeps the author's judgment---expressed as markers in the semantics
itself---and hands input generation to a backend solver, so when the specification
changes the conformance suite is \emph{regenerated}, not re-authored by
hand---a living, executable specification that drives its own testing.

\paragraph{Why error cases}
Markers target a language's \emph{negative space}: its \emph{syntactic} errors, which are ill-formed input the grammar rejects, 
and its \emph{semantic} errors---well-formed input violating the static or dynamic semantics, such as an out-of-range
value, a duplicate key, a type error, or a division by zero---the textbook
front-end/analysis boundary~\cite{pierce2002tapl}. The focus
is empirically grounded, not a convenience of what the solver finds easy.
Error-handling code is where defects concentrate and testing is thinnest: a
majority of \emph{catastrophic} distributed-systems outages trace to incorrect
handling of errors, often in code a trivial test would have
caught~\cite{yuan2014simple}, and error paths are routinely dropped in Linux
file systems~\cite{gunawi2008eio}. It is also where implementations
\emph{diverge}---no two JSON parsers agree across malformed
documents~\cite{seriot2016json}, and parser/validator discrepancies on
ill-formed input are a documented security-vulnerability
source~\cite{sassaman2013langsec,brubaker2014frankencerts}---and what
conformance suites assert normatively~\cite{tc39test262}. \rusmt{} therefore
does not fuzz values broadly; it synthesizes inputs reaching the
specification's own error and violation branches, then differentially tests how
implementations treat them~\cite{mckeeman1998differential}.

\rusmt{} collapses two of these three artifacts into one. The author writes the
executable semantics \emph{once}, in a restricted dialect of Rust whose standard
library's types denote Z3's mathematical theories. That single source is given
two interpretations: run as ordinary Rust it is the concrete conformance
\emph{oracle}; lifted to SMT-LIB it is the source of a solver query that
\emph{synthesizes} programs reaching designated points. The author marks a point
with a named marker, e.g.\ \code{Path::named("division\_by\_zero")}; \rusmt{}
asks for an input driving execution there, decodes the result into runnable
source, replays it against the oracle, and hands the witness to independent
implementations. For \toml{}, the decoded result is already a TOML document. For
\imp{}, Z3 produces an AST, so the framework renders the model back to an
\imp{} program before replay and downstream use.

When the specification changes, the suite is \emph{regenerated} by re-running
synthesis against the same markers rather than re-authored by hand. The pipeline
is intentionally ordered. Stage~1 asks Z3 to synthesize from the unconstrained
query. If Z3 returns \code{sat}, the model becomes the witness; if it returns
\code{unsat}, the marker is unreachable under the encoding and bound. Only when
Stage~1 returns no useful verdict does Stage~2 ask a model for a concrete
candidate from the emitted SMT-LIB and the solver's previous feedback. That
candidate is appended as \code{input = c}; if Z3 returns \code{sat}, the
strengthening lemma says the candidate also satisfies the original marker query.
Replay then checks that the concrete Rust oracle and the SMT lift agree on that
input. A model can guide search, but it cannot approve a test.

\paragraph{Research questions}
\textbf{RQ1:} Can a direct-style Rust interpreter, written in a small
SMT-denoting DSL, serve simultaneously as conformance oracle and marker-directed
test generator? \textbf{RQ2:} What discipline is sufficient for accepting a
generated test without human inspection? \textbf{RQ3:} When Z3 alone fails on a
realistic parser, how much coverage does the model-to-Z3 co-solving stage
recover? \textbf{RQ4:} Does the Z3 acceptance step reject model proposals often
enough to matter? \textbf{RQ5:} Run against independent implementations, what
accept/reject divergences does the generated suite expose?

\paragraph{Contributions}
This is a tool contribution: not a new solver, and not a completeness theorem.
\textbf{(1)~A Rust-native framework anyone can reuse.} You write an ordinary
recursive interpreter over the DSL vocabulary and get two artifacts from one
source: compiled by \code{rustc} it is the conformance \emph{oracle}, and lowered
to SMT-LIB it is a marker-directed \emph{test generator}---no separate grammar, no
hand-written SMT (\S\ref{sec:dsl}--\ref{sec:pipeline}). \textbf{(2)~A concrete
Z3-first, model-second pipeline.} Z3 attempts every marker first; only a timeout
or \code{unknown} invokes the model, and the model sees emitted SMT-LIB plus Z3
feedback rather than the Rust source. \textbf{(3)~Acceptance that needs no human
inspection.} A witness is kept only when Z3 shows the lifted semantics reach the
\emph{named} marker on that input, and replay checks the SMT lift against the
concrete Rust oracle (\S\ref{sec:soundness}). \textbf{(4)~Measured conformance
suites for \toml{}~1.1.0 and \imp{}.} The run covers \certcount{}/\markers{}
\toml{} markers and \impmarkers{}/\impmarkers{} \imp{} markers, with all misses
named (\S\ref{sec:eval}). \textbf{(5)~A downstream implementation comparison.}
The generated \toml{} suite is run against four real parsers and exposes
\diffcount{} accept/reject divergences. \textbf{(6)~A replayable artifact} with
pinned tool versions, generated queries, cached model transcripts, and the exact
suite used for the differential run.

% =====================================================================
\section{Background}
\label{sec:bg}

\paragraph{SMT as the reasoning backend}
\rusmt{} uses SMT to ask one operational question: \emph{does there exist an
input such that this executable semantics reaches marker \(d\)?}  The DSL
therefore maps Rust expressions to a fixed vocabulary of Z3
theories~\cite{demoura2008z3}: booleans, integers/reals, bit-vectors, strings,
sequences, sets, arrays, algebraic datatypes, and recursive functions encoded in
SMT-LIB~2.6~\cite{barrett2017smtlib}.  For each marker the backend emits a
formula of the form
\[
  \exists x.\; \mathit{entry}(x) = \mathit{Err}(\mathit{marker\_id}(d)),
\]
where \(x\) is the object-language program or input text.  On \code{sat}, Z3
returns a model for \(x\); on \code{unsat}, no such \(x\) exists \emph{for that
encoding and bound}; on \code{unknown} or timeout, the framework learns only
that this backend did not finish.

The scientific and engineering issue is that the generated formulas are outside
a decidable quantifier-free fragment.  A recursive interpreter becomes
\code{define-funs-rec}; Z3 reasons about such definitions by unfolding and
heuristics.  Parsers additionally combine strings/sequences, arithmetic,
bit-vectors, recursive datatypes, and occasional quantifiers handled by
incomplete instantiation.  Thus \rusmt{} treats Z3 as a sound decision procedure
for returned verdicts, not as a complete search oracle: \code{sat} candidates
must still replay through the Rust semantics, and \code{unknown}/timeout costs
only coverage.  The same split explains why proposer recovery is useful.  Searching
for an unconstrained input may be intractable, while validating a fully
determined candidate can be easy because the backend macro-inlines the concrete
sequence and Z3 constant-folds along one path.

\paragraph{Executable semantics and synthesis}
A \emph{semantics} fixes the meaning of a language; it is \emph{executable} when
that definition is itself a runnable artifact---a reference interpreter written
\emph{to be} the specification, not merely any implementation that happens to run
programs (a production compiler runs the language without being its definition).
Frameworks like K~\cite{rosu2010k}, PLT Redex~\cite{felleisen2009redex}, and
Ott/Lem~\cite{sewell2007ott,mulligan2014lem} write that definition as
rewrite/inference rules and derive a tool suite from them. \rusmt{} instead
takes the semantics to be an \emph{ordinary recursive function}: the same source
executes concretely as the oracle and compiles symbolically as the search space.
Compared with fuzzers such as Csmith~\cite{yang2011csmith}, the target is not
coverage in general but a named specification obligation.  Compared with SyGuS
and SemGuS~\cite{alur2013sygus,kim2021semgus}, the user does not author a
separate grammar or CHC semantics; the unknown is an input to the interpreter the
user already wrote.

% =====================================================================
\section{The \rusmt{} DSL}
\label{sec:dsl}

A \rusmt{} specification is an ordinary Rust program, written once and read two
ways: compiled by \code{rustc} it \emph{runs} as the reference interpreter;
lifted by the transpiler it \emph{denotes} an SMT formula. The design principle
is that these two readings must agree by construction, so every construct
admitted into the DSL is one with a mathematical meaning the solver can express.
That single requirement dictates the whole subset: \rusmt{} is a deliberately
\emph{functional, total} fragment of Rust.

\paragraph{A functional subset, and why each restriction earns its place}
The DSL keeps immutable \code{let} bindings, \code{if} with a mandatory
\code{else}, \code{struct}/\code{enum} with \code{match}, and (mutually)
recursive calls; it forbids loops, mutation, and references. These are not
arbitrary---each exclusion buys a property the lift depends on. \emph{No mutation
or references} means a value never aliases or changes under another name, so a
\code{let} is mathematical substitution and the translation is referentially
transparent; this is the same discipline that lets a big-step operational
semantics be read as a system of equations. \emph{No loops} leaves recursion and
(bounded) quantifiers as the sole means of iteration, matching how a solver
expresses repetition---as recursive definitions and quantified formulas, not as
mutable state evolving over time. The payoff is that the interpreter reads like
the inference rules it formalizes: Winskel's big-step rules for \imp{}
(\S\ref{sec:imp}) transcribe almost line for line, and there is no gap between
``the specification'' and ``the code that runs.''

\paragraph{Totality makes error cases first-class}
The subset is, crucially, \emph{total and exception-free}: no \code{panic!}, no
\code{?}, no \code{.unwrap()}, no early return. A partial function has no total
SMT reading, so any divergence from the happy path must be expressed as a
returned \emph{value}---for example a \code{ParseResult::Err} carrying a
\code{Path} (below)---rather than a thrown exception. This is exactly what the
negative-space thesis (\S\ref{sec:intro}) wants, and it comes for free: every
error condition is, by construction, an explicit branch returning an error value,
i.e.\ a place a marker can sit, on code that must exist anyway.

\paragraph{Separation of concerns}
Two layers are kept strictly apart. The \emph{standard library} (below) is a
fixed vocabulary denoting \emph{Z3's} theories and nothing else; the
\emph{interpreter} encodes a \emph{target} language's semantics on top of that
vocabulary. A target's idiosyncrasies---two's-complement wraparound, a parser's
error conditions, an uninitialized-variable rule---live in the interpreter as
explicit guards, never inside the stdlib. This keeps the per-construct soundness
obligation (\S\ref{sec:soundness}) small, fixed, and language-independent: it is
checked once for the stdlib, not re-argued per case study.

\paragraph{Mechanical enforcement}
Two attribute macros enforce the subset at compile time, so an out-of-subset
specification fails to build rather than mistranslating silently.
\code{\#[smt\_type]} admits only \code{struct}s and \code{enum}s whose every
value is \code{Copy} (so nothing aliases); \code{\#[smt\_fn]} rejects
\code{const}/\code{async}/\code{unsafe} and variadic functions, method receivers,
and any body beyond \code{let} bindings and a trailing expression.

Listing~\ref{lst:div} shows the whole discipline in one \imp{} fragment---the
division guard of the arithmetic evaluator. It is ordinary recursive Rust; the
only synthesis-specific token is \code{Path::named}, sitting on the error branch
a total evaluator must have anyway. The primitive \code{bv\_div} already agrees
with Z3 on every nonzero divisor, so the interpreter guards only the diverging
case (a zero divisor) and falls through to \code{bv\_div} otherwise---the
guard/if-then-else pattern of \S\ref{sec:soundness}. From this single source
\rusmt{} synthesizes \code{A := (0 / 0)} (\S\ref{sec:eval}); replayed through the
same \code{eval\_aexp}, it reaches the marker.

\begin{lstlisting}[style=rust,caption={The division guard of \imp{}'s evaluator
(operands already evaluated to \code{new\_l}, \code{new\_r}; \code{Err}
propagation elided).},label={lst:div}]
if *new_r.eq(I64::from(0)) {
    ArithmeticResult::Err(Path::named("division_by_zero"))
} else {
    ArithmeticResult::Ok(new_l.bv_div(new_r))
}
\end{lstlisting}
\paragraph{Standard library: the theory vocabulary}
The DSL's types denote Z3 sorts: \code{Boolean}, unbounded
\code{Integer}/\code{Real} (big integers/rationals), the bit-vectors
\code{I32}/\code{I64}/\code{U32}/\code{U64} (the signed/unsigned split is a
DSL-API reading of the same Z3 sort), IEEE \code{F32}/\code{F64}, \code{String}
(the Strings theory), \code{Seq<T>}, \code{Set<T>}, and \code{Array<K,V>}. As
established above, these denote \emph{Z3's} theories, not Rust's native semantics
and not any target language's. Two helper types are frontend-only. \code{Cloak<T>}
is a \code{Box}-like indirection that lets an author write self-referential ADTs
(e.g.\ \code{Add(Cloak<Aexp>, Cloak<Aexp>)}) despite Rust's sized-type rule; the
IR \emph{erases} it entirely---every \code{Cloak<T>} field lowers to \code{T} and
\code{shield}/\code{reveal} to identity---since Z3 supports recursive datatypes
natively (the one proviso is well-foundedness: a \code{Cloak<T>} must co-occur
with a non-recursive variant, so the lowered datatype is inhabited). \code{Path},
the marker type, is described next.

\paragraph{Markers via \code{Path}}
A marker flags a program point the author wants covered---characteristically an
\emph{error or specification-violation} branch (the ``negative space'' of
\S\ref{sec:intro}); markers are \emph{syntactic} or \emph{semantic}
(the taxonomy above), and the case studies span both. The value a marker
carries is a \code{Path}, and its \emph{identity} is what lets the marker be
matched across the two readings of the specification. An id drawn from a global
counter would be order-dependent, and so unmatchable against a replayed run;
instead \code{Path::named("\(d\)")} derives its id as a fixed hash of the
name~\(d\), and the \emph{same} pure function serves transpiler and evaluator, so
the id the query asserts and the id a re-executed witness carries are
\emph{identical by construction}. In the SMT backend a \code{Path} is represented
by that id as a plain \code{Integer} (with a sentinel for ``no marker reached'').
The choice is deliberate: an integer equality remains decidable through the deep
recursive unfolding of a lifted interpreter, whereas a set-membership test would
force the solver to construct and probe a symbolic collection it cannot evaluate
on those paths---an encoding decision we return to in \S\ref{sec:pipeline}. Each
named marker is thus one synthesis target (\S\ref{sec:synthesis}); named markers
are what make per-target acceptance marker-specific (\S\ref{sec:soundness}), and
the case studies we replay use them.

\paragraph{Quantifiers}
\code{forall!}, \code{exists!}, and \code{choose!} exist in a \emph{bounded}
form over a collection (\code{forall!(x in xs => p)}). \code{choose!} is
Hilbert choice: concretely the first satisfying witness, in SMT a Skolem
function under the standard choose axiom---sound because nothing may depend on
\emph{which} witness is chosen, only that it satisfies the predicate.

% =====================================================================
\section{The compilation pipeline}
\label{sec:pipeline}

The same annotated Rust source is compiled in two modes.  In oracle mode,
\code{rustc} executes the interpreter on concrete candidate inputs.  In solver
mode, the transpiler parses the restricted Rust subset, lowers it to a
sort-checked IR by Hindley--Milner unification over SMT sorts, monomorphizes
generic functions on demand, records each \code{Path::named} marker, and emits a
per-marker query.  Listing~\ref{lst:pipeline} shows the backend loop.  Every
query is inspectable SMT-LIB2, which is what makes each reported result
re-runnable with a bare \code{z3 -smt2} on a committed file.

\begin{lstlisting}[style=rust,caption={Backend loop for one entry function.
Stage~1 uses the emitted query unchanged. Stage~2 runs only after timeout or
\code{unknown}; it pins a proposed candidate as one extra equality and asks Z3 to
decide the strengthened query.},label={lst:pipeline}]
for each named marker d in entry:
    Q = lower_to_smtlib(entry, marker_id(d), k)
    r = run("z3 -smt2 -t:<budget_ms>", Q.main)
    if r == sat:
        c = decode_model(r.model)
        if replay(c, d) == ReachedTarget:
            accept(c)
        else:
            report_replay_rejection(d, c)
        continue

    if r == unsat:
        report_unreachable_under_encoding_or_bound(d)
        continue

    // Stage 2: model proposes, Z3 decides.
    for round in 1..R:
        c = model(reads: Q.main, d, history)
        v = run_z3(Q.main + assert(input_0 == c), budget)
        if v == sat && replay(c, d) == ReachedTarget:
            accept(c)
            break
        feedback(v, replay(c, d))
\end{lstlisting}

\paragraph{Backend and budgets}
The backend renders inspectable \code{.smt2} files from the IR and drives
\code{z3} as a subprocess. Stage~1 uses the unconstrained marker query exactly
as emitted. Stage~2 is reached only after Stage~1 times out or reports
\code{unknown}; the model receives the SMT-LIB text, the target marker name, and
the previous rejected candidates and Z3 verdicts. It runs in an empty working
directory with file and shell tools disabled, so it cannot inspect the Rust
semantics. The proposed input is pinned by adding one equality to the original
query; \code{unsat} means only that this candidate does not reach this marker,
while \code{sat} gives a marker-specific witness. The reported \toml{} run used
the candidate-acceptance mode with \code{RUSMT\_ROUNDS=9},
\code{RUSMT\_WITNESSES=1}, a 15\,s Z3 budget for pinned candidates, and cached
model transcripts. \imp{} uses native recursion (\code{k=0}); \toml{} also uses
\code{k=0} because the free input is the document text.

\paragraph{Engineering for Z3}
Several encodings are deliberately solver-facing.  A marker is an \code{Integer}
id, not a symbolic set membership test, because equality survives deep recursive
unfolding.  Recursive maps such as \toml{} tables are encoded as array-free
inductive association lists; an \code{Array} whose range is
a recursive datatype made even concrete-candidate validation return
\code{unknown} in early \toml{} experiments.  Recursion is offered either as
native \code{define-funs-rec} or as fixed-depth unrolling because neither
dominates across the case studies.  The backend parses the leading Z3
status line before trusting process exit status, since Z3 can exit nonzero after
benign model-output failures.  A \code{sat} model is then rendered back to
object-language source by re-sharing \code{let} terms, unmangling constructors,
canonicalizing solver-chosen names, and decoding bit-vectors as
two's-complement integers.
% =====================================================================
\section{Soundness of the lift}
\label{sec:soundness}

\paragraph{The trusted computing base}
The trusted base is the \rusmt{} specification together with the transpiler (the
lift to SMT) and Z3. We \emph{assume} the prose-to-DSL transcription correct
(trusted, not verified), and we trust the transpiler and solver \emph{because}
the lift is faithful---the per-construct contract below, \emph{checked} rather
than asserted---so trusting Z3's verdict reduces to that one obligation. The
concrete parser and pretty-printer are also trusted; a bug in any is a bug in the
results. Replay is not what discharges this trust: it is a per-candidate
cross-check that the symbolic and concrete readings agree, catching a
faithfulness defect on any candidate it touches as a spurious rejection.
Per-target acceptance additionally assumes \code{marker\_id} is injective over
the finite, author-declared set of marker names---as a $63$-bit hash, a checkable
side condition trivially met by the names we use.

\paragraph{The per-construct contract}
Soundness rests on a contract for every standard-library primitive:
\begin{definition}[Per-construct soundness contract]\label{def:contract}
For every stdlib function $f$ and every concrete input $x$,
\[
  \mathit{rust\_impl}(f, x) \;=\; \mathit{z3\_eval}\bigl(\mathit{z3\_formula}(f),\, x\bigr).
\]
\end{definition}
The stdlib denotes Z3's theories; bridging it to a given target language is the
interpreter author's job.

\paragraph{The guard / if-then-else pattern}
A primitive matching the target language is used directly; one that diverges only
on specific inputs is \emph{guarded} on those inputs, falling through to the
(Z3-correct) primitive otherwise. Listing~\ref{lst:div} has this shape:
\code{bv\_div} agrees with Z3 everywhere, so the only special case is a zero
divisor, which the guard turns into a marker. Because concrete Rust and Z3 take
the \emph{same} branch on the same input, Definition~\ref{def:contract}'s equality
survives the conditional---and the same discipline appears inside the stdlib (e.g.\
a width-or-more shift returns $0$ to match Z3's \code{bvshl}).

\paragraph{How the contract is validated}
Validation is layered: (i)~each stdlib method's doc-comment states the exact Z3
formula it must match; (ii)~example-based unit tests pin the Rust side to Z3's
documented result and have caught real divergences; (iii)~for the error-prone
\emph{bit-vector} families (signed vs.\ unsigned division/comparison; over-width
shifts; \code{i64::MIN\,/\,-1}) we discharge the contract \emph{empirically} with a
randomized differential harness comparing $\mathit{rust\_impl}$ against the
\emph{backend-emitted} formula (\S\ref{sec:eval}); the cross-backend \code{both}
mode is an extra whole-query check that catches pipeline regressions.

\paragraph{Single free input; fixed context}
A load-bearing point governs \emph{which} inputs are free in the synthesis query.
\begin{principle}[Free/fixed agreement]\label{prin:free}
The set of inputs treated as free by the SMT query must coincide with the set of
inputs left free by the concrete oracle; everything else must be fixed to the
same value on both sides.
\end{principle}
The \imp{} entry point \code{eval\_com(c: Com)} fixes the empty store
internally, so the program is the single free input; were the store left
free, Z3 could choose a starting state the concrete run never uses---a
``witness'' that does not replay. \toml{}'s \code{parse\_toml} needs no wrapper.

\paragraph{Acceptance by Z3, cross-checked by replay}
\begin{observation}[Marker-specific acceptance]
A candidate is accepted for target \(d\) only if Z3 decides that the lifted
semantics reach the \emph{named} marker \(d\), and the same candidate,
re-executed through the concrete interpreter, fires the \emph{same} \(d\)---not
merely some marker.
\end{observation}
Z3 decides the lifted reachability query; replay is an \emph{independent}
cross-check of the identical fact through the concrete reading of the spec. It
re-parses and re-executes the decoded input under a wall-clock bound and checks
membership of the \emph{same} \code{marker\_id($d$)} the query asserted---so a
candidate firing a \emph{different} marker is rejected. The cross-check is
load-bearing for two reasons. First, it \emph{guards against} a violation of the
lift obligation
(Definition~\ref{def:contract}) on whatever constructs the candidate exercises:
were the transpilation unfaithful on one, Z3 could validate a candidate the
concrete semantics do not, and replay catches exactly that disagreement (the
known stdlib divergence \code{Array::contains\_key} on sentinel values thus causes
only a \emph{spurious rejection}, never a false accept). Second, quantified or
depth-bounded queries can return a \emph{candidate} model (Z3 reports \code{sat}
with \code{(:reason-unknown "")}) that does not genuinely reach the marker;
replay rejects these too. Identifier and bit-vector decoding do not affect which
marker fires.

\paragraph{Completeness: not claimed}
We do \emph{not} claim completeness in general. Even when an input reaching a
marker exists, Z3 may return \code{unknown} or time out: the encoded formula can
fall outside what the solver decides within its resource limits. At best we have
\emph{bounded} completeness---relative to a structural/fuel bound
(\S\ref{sec:synthesis})---and even that is contingent on the solver.

% =====================================================================
\section{Bounded synthesis}
\label{sec:synthesis}

\paragraph{Markers to queries}
For a target marker with id $i$ ($i=\mathit{marker\_id}(d)$ for a named marker
\(d\)), the query declares one constant per top-level parameter (\code{input\_0},
\dots), applies the entry function, and asserts that the returned error value's
carried \code{Path} \emph{equals} $i$ (an integer equality, guarded by the enum
constructor test that selects the error variant), then
\code{(check-sat)}/\code{(get-model)}.

\paragraph{Bounded recursion by unrolling}
A recursive interpreter is not a finite formula. A \code{k=N} flag selects
between two routes. With \code{N=0} (default) recursive strongly-connected
components (SCCs) are emitted as Z3 \code{define-funs-rec} and the solver handles
recursion natively. With \code{N\,$\geq$\,1} every recursive SCC is unrolled into
$N{+}1$ depth-indexed copies $f_0,\dots,f_N$: $f_d$ calls $f_{d-1}$ for in-SCC
callees, $f_0$ is a terminator, and external callers route through $f_N$. A
\code{sat} model under depth-$N$ unrolling is a witness reachable in $\leq N$
recursive steps---the sense in which completeness is \emph{bounded}, as in bounded
model checking~\cite{biere1999bmc}.

\paragraph{A non-error terminator makes unrolling productive}
The terminator $f_0$ matters. When the return type has no nullary variant, the
backend falls back to an \emph{unconstrained} null sentinel, which the solver may
instantiate as the very error value a marker assertion demands, so a depth-$N$
\code{sat} can be satisfied \emph{through the cutoff} rather than a genuine path.
Replay rejects such candidates, so soundness is never at risk, but the budget is
wasted. The remedy is a nullary \emph{bottom} variant never produced by the
concrete semantics: the cutoff can no longer discharge an \code{Err}-marker
assertion, forcing a real marker-firing path within depth~$N$.

\paragraph{Model decoding and the loop}
A \code{sat} model is decoded to runnable source by a printer that extracts
\code{input\_0}, inlines \code{let}-shared subterms, maps mangled constructors
back to AST nodes, canonicalizes solver-chosen identifiers (sound: which marker
fires is spelling-independent), and decodes bit-vector literals as
two's-complement signed integers. This rendering step is visible in \imp{}, where
Z3 returns a \code{Com} AST that must be printed as an \imp{} program. In
\toml{}, the top-level input is already a sequence of Unicode code points, so
the decoded model is the TOML document itself. The loop per target---build,
solve, decode, replay, accept only if the marker fires---then writes the accepted
witnesses as a conformance suite.

% =====================================================================
\section{Case studies}
\label{sec:cases}

\rusmt{} is meant for authors of new executable semantics, so the case studies
are selected to stress two distinct regimes rather than to maximize benchmark
count. \imp{} is the small dynamic-semantics baseline: Z3 synthesizes an AST,
\rusmt{} renders it back to source, and replay closes the loop. \toml{} tests a
realistic parser with both syntactic and semantic errors: the input is already a
document, solver-only search fails, and the co-solving stage produces the suite
used for downstream conformance testing.

\subsection{\imp{}: closing the loop end-to-end}
\label{sec:imp}

\imp{} is the canonical small imperative ``while'' language of
Winskel~\cite{winskel1993}, chosen because its big-step semantics is familiar and
recursive but small enough that native Z3 recursion should succeed if the
pipeline is sound.  Its grammar has arithmetic and boolean expressions,
assignment, sequencing, conditionals, and while loops.  We add two dynamic error
branches to make conformance obligations explicit: division by zero and reading
an uninitialized variable from the initially empty store.  The target programs
Z3 produces are correspondingly interpretable: \code{A := (0 / 0)} for
\code{division\_by\_zero} and \code{A := v0} for
\code{uninitialized\_variable}.  This study answers whether a direct-style
interpreter can be lifted, solved, decoded, and replayed without bounded
unrolling or proposer help.

\subsection{\toml{}: a parser suite}
\label{sec:tomlcase}

\toml{}~1.1.0~\cite{tomlspec} is chosen because it is large enough to defeat the
solver while still having a precise public specification and several independent
implementations. The executable parser has a single free input, a sequence of
Unicode code points, and $151$ parser functions covering documents, key-value
pairs, tables, arrays, inline tables, booleans, integers, floats, strings, and
datetimes. Its \markers{} named markers cover syntactic errors (malformed keys,
unterminated strings, invalid escapes, missing array/table delimiters) and
semantic violations (integer/float overflow, out-of-range date/time fields,
duplicate keys, redefined tables).

This is the case study that motivates the paper's method. Free symbolic search
returns nothing for every marker, and the failure is worth stating precisely
because it dictates the remedy: given a symbolic input Z3 does not report
\code{unknown} after exhausting a budget, it fails to produce a verdict at all
within $300$\,s, and on partially-symbolic variants of the same query it
terminates abnormally. Given a \emph{determined} input the same query decides in
well under a second. The asymmetry is not about search space size but about
macro expansion: the transpiled parsers are \code{define-fun} definitions that Z3
substitutes during preprocessing, before it can know which branch is taken, so a
concrete input lets the simplifier fold each conditional as it expands while a
symbolic one leaves both arms of every conditional alive.

The recovery route follows from that asymmetry. The model supplies a determined
input from the SMT-LIB query and feedback; Z3 decides that pinned input against
the original marker query. The accepted witnesses are small and readable:
\code{a=1979-13-01} for the month-range marker,
\code{[a]\textbackslash n[a]} for the duplicate-table marker, and \code{a=[} for
an unterminated array. These documents become the conformance suite of
\S\ref{sec:eval}. We do not claim a verified complete \toml{} formalisation.

\section{Evaluation and discussion}
\label{sec:eval}

\paragraph{Evaluation questions}
The evaluation follows the RQs from \S\ref{sec:intro}: whether one executable
semantics can be both oracle and synthesis source (RQ1), whether Z3 acceptance
plus replay avoids manual inspection (RQ2), how much the co-solving stage
recovers on \toml{} (RQ3), whether Z3 rejects model proposals often enough to
matter (RQ4), and what the generated suite finds against real implementations
(RQ5).

\paragraph{Setup}
We measured on an Apple~M1 running macOS~26.5 with Z3~4.15.4, driven as a
subprocess over SMT-LIB2. The \imp{} suite was generated by
\code{cargo run -p rusmt-smt-derive -- imp eval\_com}; both markers solved in
Stage~1 under native recursion (\code{k=0}). The \toml{} suite was generated by
\code{cargo run -p rusmt-smt-derive -- toml parse\_toml} through the resumable
sweep driver with six independent workers. The \toml{} co-solving run used
\code{gpt-5.4-mini}, \code{RUSMT\_ROUNDS=9}, \code{RUSMT\_WITNESSES=1}, and a
15\,s Z3 budget for each pinned candidate. The proposer ran in an empty
directory with file and shell tools disabled, so its input was the emitted
SMT-LIB excerpt, the marker name, and Z3 feedback from previous rounds. We use
cached transcripts for replay; exact re-generation may vary because the model is
nondeterministic, so the artifact records both the generated suite and the prompt
cache used for the reported run.

\begin{table}[t]
\centering
\caption{Main synthesis and conformance results. Stage~2 counts only candidates
that Z3 decided against the original marker query with one added input equality.}
\label{tab:results}
\scriptsize
\setlength{\tabcolsep}{3pt}
\begin{tabular}{@{}p{0.31\columnwidth}p{0.25\columnwidth}p{0.35\columnwidth}@{}}
\toprule
Check & Setting & Result \\
\midrule
Stdlib contract & BV, String/Seq, edge & $2715+5317+216$ Z3-backed checks, $0$ disagreements \\
\imp{} Stage~1 & $k=0$ & \impmarkers{}/\impmarkers{} markers; ASTs rendered to \imp{} programs \\
\toml{} Stage~1 & $k=0$ & $0/\markers{}$ markers at affordable budgets \\
\toml{} Stage~2 & \code{gpt-5.4-mini}, 9 rounds & \certcount{}/\markers{} markers; \misscount{} misses \\
Proposal audit & \proposalcount{} decided/undecided proposals & \certcount{} accepted, \rejectcount{} Z3-rejected, \undecidedcount{} undecided; 65 duplicate proposals skipped before Z3 \\
Conformance diff & 131 tests $\times$ 4 parsers & \diffcount{} accept/reject divergences \\
\bottomrule
\end{tabular}
\end{table}

\paragraph{Contract-check results}
Before the synthesis results, we confirm the lift itself: we ran the differential
harness of \S\ref{sec:soundness} on the $64$-bit bit-vector families against
Z3~4.15.4: $13$ signed (\code{I64}) operations and the $6$ unsigned
(\code{U64}) operations whose formula differs from the signed one, comparing
$\mathit{rust\_impl}$ against the backend-emitted formula on directed edge cases
and seeded random inputs. The harness ran $2715$ operation-input pairs with
\emph{zero} disagreements. A companion harness covers ASCII \code{String} and
Unicode \code{Seq<U32>} families with $5317$ checks, again with zero
disagreements, and a directed edge-class harness adds $216$ more. This is not a
full proof of the stdlib contract: $32$-bit vectors, reals, and floats remain on
documented-formula-plus-example-test evidence.

\paragraph{\imp{}: the loop closes in Stage~1}
Under native recursion (\code{k=0}) the solver returns a model for both \imp{}
markers and every witness replays. Both witnesses are assembled by Z3 from the
datatype constructors alone, then rendered to source:
\code{B := (0 / 0)} reaches \code{division\_by\_zero}, and
\code{v0 := v0} reaches \code{undefined\_variable} from the initially empty
store. We do not run \imp{} through the implementation-diff harness; it is a
small end-to-end demonstration of the framework, not a language ecosystem with
competing production parsers.

\paragraph{\toml{}: solver-only search}
Lifting the \toml{} parser produces SMT that Z3 accepts without frontend error,
but solver-only search over the free document input produced no witness for any
of the \markers{} markers at affordable budgets. The failure mode is sharper
than ordinary timeout: with a symbolic input, macro expansion leaves both sides
of deeply nested parser branches alive before simplification can fold them; with
a fully determined input, the same query typically folds along one path and
decides quickly. Thus Stage~2 is not replacing a successful generator. It changes
the solver's role on hard parser targets from unconstrained search to deciding a
specific candidate against the original marker query.

\paragraph{\toml{}: co-solving coverage}
The Stage~2 run produced \certcount{} accepted TOML documents out of \markers{}
markers. The \misscount{} misses are not spread uniformly. They concentrate in
three families: date/time and numeric-offset parsing (21 markers), numeric
literal edge cases, especially float overflow and malformed radix forms (16
markers), and table/key redefinition or unterminated composite constructs (8
markers). The rest are isolated boolean, array-comment, and multiline-string
cases. This distribution is useful for diagnosis: Stage~2 handles many local
syntax errors, but remains weakest where the intended marker is easily shadowed
by an earlier parse error or where the input must navigate a long chain of
numeric subparsers before reaching the labelled branch. The accepted witnesses
are the suite; we do not fill misses with hand-written tests, and the artifact
records the full marker-level ledger.

\paragraph{Why the Z3 step matters}
Of \proposalcount{} non-duplicate candidate proposals reaching the Z3 acceptance
step, only \certcount{} became witnesses. Z3 rejected \rejectcount{} as
\code{unsat}: those inputs were valid, reached another marker, or were shadowed
by an earlier error. Another \undecidedcount{} did not produce a useful Z3
verdict within budget. A model-only generator would have had no marker-exact way
to separate these from accepted witnesses. Dynamic execution alone can say
whether an input triggers some behaviour, but it cannot prove a marker
unreachable on \code{unsat}, and it does not give the monotone acceptance
property: every model of $Q \land input=c$ is a model of the original query $Q$.

\paragraph{Conformance diff}
We ran the \certcount{} generated TOML documents against four implementations:
the Rust \code{toml} crate \code{1.1.3+spec-1.1.0}, CPython
\code{tomllib} from CPython~3.14.4, BurntSushi \code{toml}~v1.6.0 for Go, and
\code{smol-toml}~1.7.1 for Node. Table~\ref{tab:diff} summarizes agreement with
the \rusmt{} oracle. The suite found \diffcount{} accept/reject divergences over
$131 \times 4$ implementation runs.

\begin{table}[t]
\centering
\caption{Implementation agreement on the generated \toml{} suite.}
\label{tab:diff}
\scriptsize
\setlength{\tabcolsep}{4pt}
\begin{tabular}{@{}lcc@{}}
\toprule
Implementation & Agreement & Divergences \\
\midrule
Rust \code{toml} & $130/131$ & 1 \\
CPython \code{tomllib} & $123/131$ & 8 \\
BurntSushi \code{toml} & $128/131$ & 3 \\
Node \code{smol-toml} & $128/131$ & 3 \\
\bottomrule
\end{tabular}
\end{table}

The clearest candidate implementation bugs are syntax and integer-boundary
cases: \code{smol-toml} accepts \code{[[a]}, where the oracle and the other three
parsers require the closing array-table delimiter; BurntSushi accepts
\code{a = \{ b = 1 \}\textbackslash n a.c = 2}, where the oracle treats the inline
table as closed; and CPython accepts several integers outside the TOML
64-bit signed range, including \code{+9223372036854775808},
\code{-9223372036854775809}, \code{0x8000000000000000}, and oversized binary and
octal forms. Other divergences require care before calling them parser bugs:
several implementations accept extremely large positive or negative float
exponents that the current \rusmt{} formalisation rejects, so we report those as
oracle-policy/specification-triage cases rather than confirmed implementation
defects. The artifact keeps the full per-input diff so maintainers can be sent
minimal reports with the exact witness and oracle verdict.

\paragraph{Answers to the evaluation questions}
RQ1 is positive for both case studies: one Rust semantics is the oracle and the
source of generated tests. RQ2 is answered by marker-specific Z3 acceptance plus
replay; no witness is accepted from model text alone. RQ3 is positive but not
complete: co-solving covers \certcount{}/\markers{} TOML markers after Z3-alone
covers none. RQ4 is answered by the rejection count: \rejectcount{} proposed
inputs were rejected by Z3, so the solver is doing real work in the loop. RQ5 is
the differential result: the generated suite exposes \diffcount{} accept/reject
divergences across four TOML parsers.

% =====================================================================
\section{Related work}
\label{sec:related}

\paragraph{Solver-aided executable semantics}
Rosette~\cite{torlak2013rosette,torlak2014rosette} is the closest
solver-aided host: a symbolic VM compiles a language subset to constraints while
concrete evaluation strips away ordinary host code.  \rusmt{} is narrower and
more conformance-oriented: the user writes a reference interpreter, marks
specification obligations, and receives tests accepted only by Z3 validation plus
replay.  SemGuS~\cite{kim2021semgus} synthesizes from syntax and semantics
encoded as constrained Horn clauses; \rusmt{} instead derives formulas from
ordinary recursive functions.  Synantic~\cite{liu2024synantic} synthesizes a
semantics from a black-box interpreter, complementary to our source-to-SMT route.

\paragraph{Specification-derived testing}
Executable-semantics frameworks such as K~\cite{rosu2010k}, PLT
Redex~\cite{felleisen2009redex}, and Ott/Lem~\cite{sewell2007ott,mulligan2014lem}
derive tools from rewrite or inference-rule specifications.  \rusmt{} derives a
more focused tool from direct-style Rust: marker-directed conformance tests.
JEST~\cite{park2021jest} is closest in purpose, using a mechanized ECMAScript
spec as oracle, synthesizer, and conformance tester.  Our distinction is the
Z3/replay acceptance gate on author-named markers and the explicit
sound-but-incomplete split exposed by the \toml{} case study.  Coverage- and
grammar-driven generators such as Csmith~\cite{yang2011csmith} remain the
application domain rather than the technique: \rusmt{} derives targeted witnesses
from the executable specification.

\paragraph{Language models behind trusted checkers}
We do not claim novelty for the proposer/checker pattern.  Lemur~\cite{wu2024lemur}
uses LLM-proposed verification structure, and AquaForte~\cite{lv2026aquaforte}
has an LLM propose SMT models that the solver rechecks.  Test-generation
portfolios such as CodaMOSA~\cite{lemieux2023codamosa} and
Cottontail~\cite{tu2025cottontail} similarly call a model only when conventional
generation stalls.  \rusmt{}'s modest difference is where the checked obligation
lives: the same executable semantics is oracle and synthesis source, and the
accepted obligation is a stable named marker rather than generic coverage. The
model is useful as a generator after solver search stalls; Z3 still decides each
accepted candidate against the emitted marker query.


% =====================================================================
\section{Limitations and threats to validity}
\label{sec:limits}

\paragraph{Scope and subset}
The DSL excludes loops, mutation, and references; the price buys a clean SMT
translation, mirrors textbook big-step semantics, and is mechanically enforced.
\rusmt{} targets a specification's \emph{negative space}---its syntactic and
semantic error branches---not happy-path fuzzing (\S\ref{sec:intro}). A marker
covers only an error condition the author \emph{wrote}; a violation the
specification implies but the reference omits is not marked. The \toml{} suite is
therefore marker-complete only relative to this formalisation, not to all of TOML
1.1.0. Naming \toml{} markers does not change the lifted query, only which
targets are name-addressable.

\paragraph{The authoring burden}
The artifact is \emph{write-once, use-many}---one source is oracle, synthesis
source, and replay oracle, amortizing over a reusable specification rather than
a throwaway generator~\cite{rosu2010k,felleisen2009redex}. The burden is lowered
but not removed: a human must choose the DSL-expressible subset, place the
markers, and review and own the result. The approach thus pays off when a
reusable reference semantics is wanted anyway, and is poor economics for a
one-off test.

\paragraph{Soundness validation and the solver frontier}
The per-construct contract is \emph{checked} only for the bit-vector and
string/sequence families (\S\ref{sec:eval}), not the $32$-bit variants, reals, or
floats, where a divergent primitive would silently degrade faithfulness;
extending the harness is the most important remaining validation work. We claim
only \emph{bounded} completeness, contingent on the solver. The dominant
limitation is solver capability on deeply recursive, multi-theory queries, and we
now characterise it precisely rather than reporting it as a wall: Z3 decides a
\emph{determined} input on the \toml{} parser quickly and returns nothing useful
for a symbolic one at affordable budgets, because the transpiled parsers are
macros expanded before branch conditions can fold. That diagnosis is what makes
the co-solving stage principled rather than opportunistic---the model attacks
search, and Z3 still decides acceptance---but it does not make the solver
complete, and \misscount{} markers remain uncovered (\S\ref{sec:eval}). A natural
next step lets a stuck solver hand control back
through Z3's user-propagator API---callbacks at decision points the solver
controls, no vendored patch---so a proposer could inject solver-\emph{rechecked}
lemmas~\cite{lv2026aquaforte,wu2024lemur}, or, more ambitiously, register the
concrete reference semantics as a propagated theory so a candidate is replayed
\emph{inside} the search and a rejection returns a conflict clause rather than
merely ending a round. Notably this needs no second encoder:
\code{Z3\_solver\_from\_string} loads the backend's own emitted SMT-LIB into an
in-memory solver to which a propagator attaches directly.

\paragraph{Scope of the AI integrations}
The co-solving stage is measured, not only described: every reported recovery comes
from live model calls, and the model is a pipeline stage rather than an optional
extra. What we do \emph{not} claim is a comparison between models, an
authoring-effort measurement, or that the model's proposals are reliable in
isolation---they are not, and they do not need to be, because Z3 decides each one
against the unmodified query. The proposer runs with its file and shell tools
disabled in an empty working directory, so it cannot read the Rust semantics and
recover an answer from the source rather than from the emitted SMT-LIB; we
verified this directly. Every exchange is recorded, keyed by a content hash of
the prompt, so the reported run replays without model access. Re-running the
model from scratch is not guaranteed to regenerate the identical suite; the
scientific result is therefore the recorded run plus the generated suite, not a
claim of model determinism.

\paragraph{Threats to validity, by category}
Read against the standard empirical-software-engineering
taxonomy~\cite{wohlin2012experimentation,runeson2009casestudy}:
\emph{construct}---a marker measures an error the reference author \emph{wrote},
not the specification's true error space, and the per-construct contract is checked
only for the bit-vector and string/sequence families ($2715{+}5317{+}216$ Z3-backed
checks), so ``the lift is faithful'' is measured there and argued for reals,
floats, and $32$-bit vectors. \emph{Internal}---acceptance is by replay under
\emph{named} markers whose id coincides on both sides by construction, so a witness
is credited only when re-execution fires \emph{that} target, excluding
``some-marker'' confounds and spurious \code{k}-bounded models (\S\ref{sec:eval});
the array-free re-encoding is controlled by byte-identical concrete output across
both encodings on every test (regression-preserving, not formally proven
equivalent). \emph{External}---two case studies, one machine, one solver version,
one model, no cross-machine or cross-model claim; the frontier is exhibited by
\toml{} but argued from recursion depth $\times$ theory combination to be a
property of method$\,\times\,$solver. \emph{Conclusion}---the \toml{} co-solving
run is a single recorded model run, and Z3's heuristics make latencies
run-dependent, so we report the qualitative regime and content-address each
query; ``zero disagreement'' attests no counterexample \emph{within} the checks,
not exhaustive equivalence.

% =====================================================================
\section{Conclusion}
\label{sec:concl}

\rusmt{} lets an author write a language's reference semantics once, as an
ordinary recursive Rust program that is at the same time the conformance oracle
and the source of the tests. Z3 alone closes the loop for \imp{}; on the larger
\toml{} parser, a model proposes concrete documents from the emitted SMT-LIB and
Z3 decides whether each one reaches the named marker. The result is a generated
\toml{} conformance suite covering \certcount{}/\markers{} markers and exposing
\diffcount{} accept/reject divergences across four parsers. The model is not a
judge and the suite is not complete, but the accepted tests are marker-specific
and replayed against the concrete oracle. That is the core promise: living
executable specifications that regenerate targeted conformance suites when the
spec changes.

\bibliographystyle{ACM-Reference-Format}
\bibliography{refs}

% =====================================================================
\end{document}
